%%%%%%%%%%%%%%%%%%%%%%%%%%%%%%%%%%%%%%%%%%%%%%%%%%%%%%%%%%%%%%%%%%%%
%% I, the copyright holder of this work, release this work into the
%% public domain. This applies worldwide. In some countries this may
%% not be legally possible; if so: I grant anyone the right to use
%% this work for any purpose, without any conditions, unless such
%% conditions are required by law.
%%%%%%%%%%%%%%%%%%%%%%%%%%%%%%%%%%%%%%%%%%%%%%%%%%%%%%%%%%%%%%%%%%%%

\documentclass{beamer}
\usetheme[faculty=ped,basePath=../fibeamer,logo=../fibeamer/logo/mu/Logo-UC-3.png]{fibeamer}
\graphicspath{{../resources/}}

\usepackage{algorithm}
\usepackage{algpseudocode}

\usetheme[faculty=ped]{fibeamer}
\usepackage[utf8]{inputenc}
\usepackage[
  main=english, %% By using `czech` or `slovak` as the main locale
                %% instead of `english`, you can typeset the
                %% presentation in either Czech or Slovak,
                %% respectively.
  czech, slovak %% The additional keys allow foreign texts to be
]{babel}        %% typeset as follows:
%%
%%   \begin{otherlanguage}{czech}   ... \end{otherlanguage}
%%   \begin{otherlanguage}{slovak}  ... \end{otherlanguage}
%%
%% These macros specify information about the presentation
\title{Tuplas y listas: ejercicios} %% that will be typeset on the
\subtitle{Semestre II-2026} %% title page.
\author{Profesora: Andreina Cota Vidaurre}
%% These additional packages are used within the document:
\usepackage{ragged2e}  % `\justifying` text
\usepackage{booktabs}  % Tables
\usepackage{tabularx}
\usepackage{tikz}      % Diagrams
\usetikzlibrary{calc, shapes, backgrounds}
\usetikzlibrary{positioning, arrows.meta, shadows.blur, shapes.geometric}
\usepackage{amsmath, amssymb}
\usepackage{url}       % `\url`s
\usepackage{listings}  % Code listings
\usepackage{graphicx}
\usepackage{amsmath}
\usepackage[T1]{fontenc}
\usepackage{xcolor}
\usepackage{minted}
\usepackage{pifont}
% \usepackage{adjustbox}


% Definicion de pseudocodigo en espaniol
\lstdefinelanguage{PseudocodigoES}{
    morekeywords={INICIO,FIN,PARA,HACER,SI,ENTONCES,FIN_SI,FIN_PARA,IMPRIMIR},
    sensitive=true,
    morecomment=[l]{//},
    morestring=[b]"
}

% Definicion de pseudocodigo en ingles
\lstdefinelanguage{PseudocodeEN}{
    morekeywords={BEGIN,END,FOR,DO,IF,THEN,END_IF,END_FOR,PRINT},
    sensitive=true,
    morecomment=[l]{//},
    morestring=[b]"
}

\lstset{
    basicstyle=\ttfamily\small,
    keywordstyle=\bfseries,
    columns=fullflexible,
    keepspaces=true,
    showstringspaces=false,
    frame=single,
    xleftmargin=0em,
    framexleftmargin=0em,
    framesep=2pt,
    gobble=4,
    resetmargins=true,
    aboveskip=0pt,
    belowskip=0pt
}

\lstdefinestyle{pythonstyle}{
    language=Python,
    basicstyle=\ttfamily\small,
    keywordstyle=\bfseries,
    commentstyle=\itshape,
    stringstyle=\ttfamily,
    showstringspaces=false,
    columns=fullflexible,
    keepspaces=true,
    frame=single,
    breaklines=true,
    xleftmargin=0em,
    framexleftmargin=0em,
    framesep=2pt,
    aboveskip=0pt,
    belowskip=0pt,
    gobble=4,
    resetmargins=true,
    emph={print,input,int,float,str,bool,len,range},
    emphstyle=\bfseries
}

\lstset{style=pythonstyle}

% Estilos del diagrama
\tikzstyle{iniciofin} = [
    ellipse,
    minimum width=2.5cm,
    minimum height=1cm,
    text centered,
    draw=black,
    fill=blue!25
]

\tikzstyle{proceso} = [
    rectangle,
    rounded corners,
    minimum width=3cm,
    minimum height=1cm,
    text centered,
    draw=black,
    fill=cyan!20
]

\tikzstyle{decision} = [
    diamond,
    aspect=2,
    text centered,
    draw=black,
    fill=yellow!35
]

\tikzstyle{flecha} = [
    thick,
    ->,
    >=Stealth
]

% \definecolor{green_python}{HTML}{52A736}

% \newcolumntype{Y}{>{\raggedright\arraybackslash}X}

\frenchspacing
\begin{document}
  \shorthandoff{-}
  \frame[c]{\maketitle}

% \begin{darkframes}

    \begin{frame}[fragile, t]{Nivel 1: Rellenar el código}
        \begin{enumerate}
            \item Se tiene una lista de clases de un curso. Completa el código para agregar 
            \mintinline{python}{"programacion"} al final de la lista y luego imprimirla.
            \begin{minted}[xleftmargin=-2cm, frame=leftline, fontsize=\small]{python} 
            clases = ["matematicas", "fisica", "quimica"]
            # Agrega "programacion" al final de la lista
            clases.______("programacion")
            print(clases)
            # ['matematicas', 'fisica', 'quimica', 'programacion']
            \end{minted}
            \item Completa la tupla con la coordenada que falta y accede al segundo elemento.
            \begin{minted}[xleftmargin=-2cm, frame=leftline, fontsize=\small]{python}
            # Coordenada (latitud, longitud, altitud)
            posicion = (4.7110, _______, 2600)
            # Accede a la longitud (segundo elemento)
            longitud = posicion[_______]
            print(longitud) # -74.0721
            \end{minted}
        \end{enumerate}
    \end{frame}

    \begin{frame}[fragile, t]{Nivel 1: Rellenar el código}
        \begin{enumerate}
            \item[3.] Completa el código para conocer cuántas notas hay y cuál es la primera nota.
            \begin{minted}[xleftmargin=-2cm, frame=leftline, fontsize=\small]{python}
            notas = [85, 92, 78, 95, 60]
            
            cantidad = _______(notas)
            
            primera_nota  = notas[_______]
            
            print(cantidad, primera_nota) # 5   85
            \end{minted}
        \end{enumerate}
    \end{frame}

    \begin{frame}[fragile, t]{Nivel 2: Uso de funciones}
        \begin{enumerate}
            \item[4.] Tienes las temperaturas registradas durante una semana en una localidad. 
            Usa funciones para encontrar la máxima, la mínima y el promedio de las temperaturas.
            \begin{minted}[xleftmargin=-2cm, frame=leftline, fontsize=\small]{python}
            temperaturas = [18, 22, 15, 25, 19, 13, 21]
            
            maxima  = ___(temperaturas)
            minima  = ___(temperaturas)
            promedio = sum(temperaturas) / ___(temperaturas)
            
            print(f"Máxima temp: {maxima}°C")
            print(f"Mínima temp: {minima}°C")
            print(f"Promedio temp: {promedio:.1f}°C")
            \end{minted}
            ¿Cómo ordenarías las temperaturas de mayor a menor?
        \end{enumerate}
    \end{frame}

    \begin{frame}[fragile, t]{Nivel 2: Uso de funciones}
        \begin{enumerate}
            \item[5.] Dos puntos en un mapa se almacenan como tuplas. Calcula la distancia euclidiana entre ellos. 
            La fórmula para calcular la distancia es: $\sqrt{(x_2-x_1)^{2} + (y_2-y_1)^{2}}$
            \begin{minted}[xleftmargin=-2cm, frame=leftline, fontsize=\small]{python}
            import math

            punto_a = (3, 4)
            punto_b = (7, 1)
            
            distancia = ______________________________
            print(f"Distancia: {distancia:.2f} unidades")
            \end{minted}
        \end{enumerate}
    \end{frame}

    \begin{frame}[fragile, t]{Nivel 2: Uso de funciones}
        \begin{enumerate}
            \item[6.] Un dron se encuentra en la siguiente coordenada: \mintinline{python}{posicion = (15,30)}. 
            Completa el código para mostrar las coordenadas X y Y del dron.
            \begin{minted}[xleftmargin=-2cm, frame=leftline, fontsize=\small]{python}
            print("X:", posicion[___])
            print("Y:", posicion[___])
            \end{minted}
            \item[7.] Se tiene una lista de productos comprados en un supermercado:
            \begin{minted}[xleftmargin=-2cm, frame=leftline, fontsize=\small]{python}
            productos = ["Pan","Leche","Huevos","Pan","Frutas"]
            \end{minted}
            \begin{itemize}
                \item Mostrar cuántos productos existen en la lista
                \item Mostrar cuantas veces aparece "Pan" en la lista
            \end{itemize}
        \end{enumerate}
    \end{frame}

    \begin{frame}[fragile, t]{Nivel 2: Uso de funciones}
        \begin{enumerate}
            \item[8.] Se tiene la lista de estudiantes de un curso.
            \begin{minted}[xleftmargin=-2cm, frame=leftline, fontsize=\small]{python}
            estudiantes = ["Juan","Pedro","Luis", "Maria", "Javiera"]
            \end{minted}
            \begin{itemize}
                \item Eliminar de la lista al estudiante "Pedro"
                \item Agregar en su lugar un nuevo estudiante llamado "Carlos"
                \item Mostrar la lista actualizada
                \item Mostrar la cantidad total de estudiantes
            \end{itemize}
            \item[9.] Se tiene una lista con las notas de un examen:
            \begin{minted}[xleftmargin=-2cm, frame=leftline, fontsize=\small]{python}
            notas=[7,6,4,5,6,7]
            \end{minted}
            \begin{itemize}
                \item Contar cuántas veces aparece la nota 6
                \item Mostrar en qué posición aparece la primera vez la nota 6
            \end{itemize}
        \end{enumerate}
    \end{frame}

    \begin{frame}[fragile, t]{Nivel 2: Uso de funciones}
        \begin{enumerate}
            \item[10.] Tienes los precios de varios productos. Calcula el total de la compra y verifica si alcanza el presupuesto.
            \begin{minted}[xleftmargin=-2cm, frame=leftline, fontsize=\small]{python}
            precios = [1500, 3200, 800, 4500, 950]
            presupuesto = 12000
            
            total = sum(precios)
            print(f"Total: ${total}")
            print(f"Presupuesto: ${presupuesto}")
            # completar el codigo
            \end{minted}
        \end{enumerate}
    \end{frame}

    \begin{frame}[fragile, t]{Nivel 3: Condicionales}
        \begin{enumerate}
            \item[11.] Si la nota es mayor o igual a 4, el estudiante aprueba; si no, reprueba.
            Completa el código para imprimir el estado de aprobación de cada estudiante.
            \begin{minted}[xleftmargin=-2cm, frame=leftline, fontsize=\small]{python}
            notas = [7, 3.7, 6, 3, 6.5, 4]

            for nota in notas:
                # completar el codigo
                pass
            \end{minted}
        \end{enumerate}
    \end{frame}

    \begin{frame}[fragile, t]{Nivel 3: Condicionales}
        \begin{enumerate}
            \item[12.] Un sistema de acceso verifica si un soldado está en la lista autorizada. 
            Si el nombre se encuentra en la lista, imprimir: \textcolor{blue}{Apellido}: Acceso concedido; 
            si no se encuentra, imprimir: \textcolor{blue}{Apellido}: Acceso denegado.
            
            \begin{minted}[xleftmargin=-2cm, frame=leftline, fontsize=\small]{python}
            autorizados = ["García", "López", "Martínez", "Rodríguez"]
            apellido = "López"
            # completar el codigo
            \end{minted}
        \end{enumerate}
    \end{frame}

    \begin{frame}[fragile, t]{Nivel 4: Ciclos}
        \begin{enumerate}
            \item[13.] Calcula cuántas raciones hay disponibles en total acumulando día a día.
            \begin{minted}[xleftmargin=-2cm, frame=leftline, fontsize=\small]{python}
            raciones_diarias = [120, 95, 110, 88, 130, 102, 115]
            acumulado = 0
            
            for i in range(len(raciones_diarias)):
                acumulado += raciones_diarias[i]
                print(f"Día {i+1}: {raciones} raciones — \
                Acumulado: {acumulado}")
            \end{minted}
        \end{enumerate}
    \end{frame}

    \begin{frame}[fragile, t]{Nivel 4: Ciclos}
        \begin{enumerate}
            \item[14.]Genera pares \mintinline{python}{(n, resultado)} para la tabla de multiplicar de un número 
            usando una lista de tuplas. Completa el código para generar la tabla de multiplicar del número 7.
            \begin{minted}[xleftmargin=-2cm, frame=leftline, fontsize=\small]{python}
            numero = 7
            tabla = []
            
            for i in range(1, 11):
                # completar el codigo
                pass
            
            for factor, resultado in tabla:
                print(f"{numero} x {factor} = {resultado}")
            \end{minted}
        \end{enumerate}
    \end{frame}

    \begin{frame}[fragile, t]{Nivel 4: Ciclos}
        \begin{enumerate}
            \item[15.] De una lista de edades de estudiantes, crea una nueva lista solo con aquellos que tienen entre 18 y 28 años.
            \begin{minted}[xleftmargin=-2cm, frame=leftline, fontsize=\small]{python}
            edades = [17, 22, 31, 19, 26, 15, 28, 35, 21]
            lista_filtrada = []
            
            for edad in edades:
                # completar el codigo
                pass
            
            print(f"Lista filtrada: {lista_filtrada}")
            print(f"Total de estudiantes en lista: {len(lista_filtrada)}")
            \end{minted}
        \end{enumerate}
    \end{frame}
    
\end{document}
